\section{Introduction}
\label{sec:intro}

All three experiments use models small enough to train or sample on a single GPU: Qwen3-0.6B, on its own and as the language half of the vision-language model, and the 3B Qwen2.5 models whose sampled solutions were supplied for the test-time study. At this scale a change in a score is easy to measure but harder to attribute, because one number combines what a change was designed to do, what it does by accident elsewhere, and how the score itself is produced, down to which part of an answer an extractor reads. Fine-tuning shows this most directly.

Training Qwen3-0.6B on worked mathematical solutions should improve its mathematics, but every update also moves weights that serve unrelated skills, which can then degrade. The methods we compare differ only in how freely the weights may move away from their starting point $\theta_0$, and they can all be written as one constrained objective,
\begin{equation}
\min_{\theta}\;\mathcal{L}_{\mathrm{SFT}}(\theta)+\lambda\,\Omega(\theta,\theta_0)\qquad\text{subject to}\qquad \theta-\theta_0\in\mathcal{S},
\label{eq:ft}
\end{equation}
where $\mathcal{L}_{\mathrm{SFT}}$ is the loss on the training responses, $\Omega$ penalises distance from the start and $\mathcal{S}$ restricts which updates are allowed. Full fine-tuning leaves every weight free. L2-SP~\citep{li2018l2sp} keeps them free but pulls them back towards $\theta_0$. LoRA~\citep{hu2022lora} allows only a low-rank correction to each weight matrix. OPLoRA~\citep{xiong2026oplora} also keeps that correction away from the matrix's strongest existing directions. The goal is to measure the gain in mathematics (adaptation) against what is kept elsewhere (retention), and to relate both to how far, and in which directions, each method moved the weights.

The second experiment connects a frozen image encoder to the same language model. A small trained connector translates the encoder's patch features into vectors that the language model reads like word embeddings, and a curriculum teaches it to use them: first describing images, then simple questions about objects, then natural questions. The difficulty is that a model can often give a plausible answer without looking, because questions carry strong hints (``Is there a\dots?'' is usually answered ``yes''). Complementary pairs~\citep{goyal2017vqav2} remove that shortcut: the same question is asked about two images with different correct answers. Swapping the two images, using an unrelated image or removing the image then shows whether answers follow what is in the picture. The goal is to establish how strongly the model's answers depend on the image, how each curriculum stage contributes, and whether likelihood-based and generated-answer measurements agree.

The third experiment leaves the weights untouched and spends computation at inference time. Sampling many solutions to a mathematics problem raises the chance that one of them is right, but only a selection rule that recognises the right one turns that chance into accuracy. We compare voting over final answers with rules that rank solutions by the model's own confidence~\citep{kang2025selfcertainty} or, for a vision model, by how much it attends to the image~\citep{usman2026visconf}, either alone or as weights inside a vote. Alongside these offline analyses, a small live experiment compares spending the same output allowance on one long answer, three short ones, or an answer followed by a self-review, and a final study asks whether sampling can stop early once the answers agree. The goal is to locate what limits the benefit of extra samples, and how much of it a cheap rule can capture.

Section~\ref{sec:setup} summarises the shared setup and deviations, Sections~\ref{sec:task1} to~\ref{sec:task3} analyse the three experiments, and Section~\ref{sec:conclusion} concludes.
